%% =====================================================================
%%  SUPPLEMENTARY INFORMATION
%%  Confidence-Guided Adaptive Reconstruction for Windowed QPE
%%  Companion to A_method.tex. Standalone document.
%%  Build:  pdflatex supplementary / bibtex supplementary / pdflatex x2
%% =====================================================================
\documentclass[preprint,11pt]{elsarticle}

\usepackage{amsmath,amssymb,amsthm}
\usepackage{booktabs}
\usepackage{graphicx}
\usepackage{xcolor}
\usepackage{array}
\usepackage[hidelinks]{hyperref}

\newcommand{\PLACE}[1]{\textcolor{blue}{\textbf{[PLACEHOLDER: #1]}}}

% evidentiary labels, identical to the main article
\newcommand{\PROVEN}{\textsc{(proven)}}
\newcommand{\HEUR}{\textsc{(heuristic)}}
\newcommand{\EMPIR}{\textsc{(empirical)}}
\newcommand{\INVAR}{\textsc{(implementation invariant)}}

\newcommand{\Cw}{C_w}
\newcommand{\dcov}{\delta_{\mathrm{cov}}}
\newcommand{\Smax}{S_{\max}}
\newcommand{\beammax}{\texttt{beam\_max}}

% S-prefixed numbering so nothing collides with the main article
\renewcommand{\thesection}{S\arabic{section}}
\renewcommand{\thetable}{S\arabic{table}}
\renewcommand{\thefigure}{S\arabic{figure}}
\renewcommand{\theequation}{S\arabic{equation}}

\journal{MethodsX}

\begin{document}

\begin{frontmatter}
\title{Supplementary Information for\\
\emph{Confidence-Guided Adaptive Reconstruction for Windowed Quantum
Phase Estimation}}
\author[inst1]{Sanjay Lakshmanan R}
\ead{rajasanjay21@gmail.com}
\address[inst1]{Department of Computer Science and Engineering,
Amrita Vishwa Vidyapeetham, Chennai, India}

\begin{abstract}
This document contains the complete mathematical derivations, the
bin-level calibration data, and the artifact map supporting the main
article. Nothing here is required to understand the method or to judge
the central claims; everything here is required to verify them
independently. No number in this document differs from the corresponding
number in the main article, and every value is a direct summary of a
released result file (Sec.~\ref{si:artifacts}). The failure-study
documentation and the extended statistical tables supporting the
end-to-end evaluation accompany the companion study~\cite{companion}
instead.

Section numbers of the form ``Sec.~S$n$'' refer to this document.
References to the main article are given by section name.
\end{abstract}
\end{frontmatter}

\tableofcontents
\vspace{1em}
% =====================================================================
\section{Complete mathematical derivations and proofs}
\label{si:proofs}

Each result is restated before its proof. The evidentiary label is the
one used in the main article: \PROVEN{} for a mathematical argument
independent of any implementation, \INVAR{} for a property of the
specific program described, \HEUR{} for a plausibility argument with no
proof. Numbering of propositions follows the main article's
\emph{Theoretical properties} subsection.

\subsection{Notation}

\begin{table}[h]
\centering\small
\caption{Notation used throughout the main article and this document. Every symbol is also defined where it first appears in the main text.}
\label{tab:notation}
\begin{tabular}{@{}ll@{}}
\toprule
Symbol & Meaning \\
\midrule
$w$, $W$ & window index; total number of windows \\
$S_w$ & shots collected for window $w$ \\
$b^{(1)}_w, b^{(2)}_w$ & top-1, top-2 observed bit patterns in window $w$ \\
$n_1, n_2$ & raw counts of $b^{(1)}_w, b^{(2)}_w$ \\
$p_w(b)$, $\hat p_w(b)$ & true and empirical probability of pattern $b$ \\
$K_w$ & number of \emph{distinct} observed patterns, $K_w\le\min(2^{\ell_w},S_w)$ \\
$\ell_w$ & window width in bits \\
$\alpha_0$ & Beta/Dirichlet smoothing constant, fixed at $0.5$ \\
$\Cw$ & posterior comparison statistic (main article, Sec.~1.4) \\
$k$, $m_w$ & candidate count: fixed (baseline) and adaptive (module B) \\
$P$, $P'$ & beam width: fixed (baseline) and adaptive (module C) \\
$\dcov$, $\varepsilon$ & coverage slack and shot-stopping targets \\
$\Smax$, $\beammax$ & hard engineering caps \\
$o$ & overlap width between adjacent windows \\
$R_w$ & resampling rounds executed by module D on window $w$ \\
\bottomrule
\end{tabular}
\end{table}

\subsection{Smoothing bound (Proposition 1)}
\PROVEN{}\quad\textit{$|\tilde p_w(b)-\hat p_w(b)|\le\alpha_0/(\alpha_0 2^{\ell_w}+S_w)$
for the Jeffreys-smoothed estimate
$\tilde p_w(b)=(\alpha_0+\mathrm{counts}_w(b))/(2^{\ell_w}\alpha_0+S_w)$.}

\begin{proof}
Writing the difference over a common denominator,
\[
\tilde p_w(b)-\hat p_w(b)=\frac{\alpha_0 S_w-\alpha_0 2^{\ell_w}\mathrm{counts}_w(b)}{S_w(\alpha_0 2^{\ell_w}+S_w)} .
\]
Bounding $0\le\mathrm{counts}_w(b)\le S_w$ gives numerator magnitude at
most $\alpha_0 S_w$; dividing by the denominator yields the claim. At
$S_w=4$, $\ell_w\le6$ this is $0.5/(0.5\cdot64+4)\approx0.015$, which is
why the implementation ranks by the raw empirical $\hat p_w$ rather than
the smoothed estimate: at every $(\ell_w,S_w)$ pair in the evaluated
corpus the perturbation is too small to reorder the ranking.
\end{proof}

\subsection{Symmetry and boundary value (Proposition 2)}
\PROVEN{}\quad\textit{$\Cw(n_1,n_2,S_w)+\Cw(n_2,n_1,S_w)=1$, and hence
$\Cw(n,n,S_w)=1/2$.}

\begin{proof}
$\Cw(n_2,n_1,S_w)=\Pr[X_2>X_1]$ under the same joint law, with only the
labels exchanged. $X_1,X_2$ are independent and continuous, so
$\Pr[X_1=X_2]=0$ and hence $\Pr[X_1>X_2]+\Pr[X_2>X_1]=1$. Setting
$n_1=n_2=n$ gives $2\Cw(n,n,S_w)=1$.
\end{proof}

\noindent
Both identities are asserted by the implementation's own self-test, which
checks $\Cw(10,10,20)=1/2$ and $\Cw(6,4,10)+\Cw(4,6,10)=1$ to within
$10^{-6}$. The proof is why those tests are expected to hold exactly
rather than approximately, up to quadrature tolerance.

\subsection{Monotonicity (Proposition 3)}
\PROVEN{}\quad\textit{$\Cw$ is non-decreasing in $n_1$ for fixed
$n_2,S_w$, and non-increasing in $n_2$ for fixed $n_1,S_w$.}

\begin{proof}
The Beta family has monotone likelihood ratio in its first shape
parameter \cite[Ch.~3]{lehmannromano2005}, so $\mathrm{Beta}(a,b)$ is
stochastically non-decreasing in $a$ for fixed $b$ and non-increasing in
$b$ for fixed $a$. Incrementing $n_1$ raises $a_1=\alpha_0+n_1$ by one
\emph{and} lowers $b_1=\alpha_0+S_w-n_1$ by one, so both effects push
$X_1$ upward in the stochastic order:
$X_1^{(n_1+1)}\succeq_{\mathrm{st}}X_1^{(n_1)}$. Stochastic dominance is
preserved under $\Pr[\,\cdot>X_2]$ for $X_2$ held fixed. The statement
for $n_2$ follows via Proposition 2.
\end{proof}

\noindent
This is the property that licenses using $\Cw$ as a ranking signal: it
guarantees the ordering it induces over windows moves in the right
direction with the evidence, irrespective of the calibration defect
measured in the main article.

\subsection{Asymptotic consistency (Proposition 4)}
\HEUR{}\quad\textit{Fix $p_1>p_2>0$ and let $n_1,n_2,S_w\to\infty$ with
$n_1/S_w\to p_1$ and $n_2/S_w\to p_2$. Then $\Cw\to1$ in probability.}

\begin{proof}[Proof sketch]
By the Bernstein--von Mises property of the Beta posterior
\cite[Ch.~10]{vandervaart1998},
$X_i\mid n_i,S_w\approx\mathcal{N}(n_i/S_w,\ \hat p_i(1-\hat p_i)/S_w)$
for large $S_w$, uniformly on compacta away from $\{0,1\}$; the
$\alpha_0$ term is $O(1/S_w)$ and vanishes. Then $X_1-X_2$ is
asymptotically $\mathcal{N}(p_1-p_2,\ O(1/S_w))$, whose variance tends to
zero while its mean is fixed at $p_1-p_2>0$. By Slutsky's theorem,
$\Pr[X_1-X_2>0]\to1$.
\end{proof}

\noindent
This is labelled a sketch rather than a proof, and the label is not
modesty: the normal approximation invoked is itself asymptotic, and no
finite-sample rate (a Berry--Esseen-type bound, for instance) is derived.
Nothing downstream in the method requires one. The implementation's
self-test checks the qualitative consequence at
$n_1{:}n_2{:}S_w=600{:}400{:}1000$, where $\Cw>0.999$.

\subsection{Module B returns the minimal covering prefix (Proposition 5)}
\PROVEN{}\quad\textit{The coverage rule returns the shortest
descending-probability prefix whose cumulative empirical mass reaches
$1-\dcov$.}

\begin{proof}
Candidates are appended in descending-probability order, and the
termination condition is checked against the cumulative mass \emph{before}
the next append. Since $\hat p_w\ge0$ and the list is sorted descending,
the partial sums are non-decreasing in $m$, so the first $m$ at which the
target is met is the unique minimal such $m$. Existence holds because
$\sum_i\hat p_w(b_{(i)})=1\ge1-\dcov$ for any $\dcov>0$. The rule
therefore returns exactly $\{b_{(1)},\dots,b_{(m_w)}\}$ with $m_w$ the
minimal index. Non-emptiness is a separate claim and also holds: the
break condition cannot fire on the first iteration, because the candidate
list is empty at that point, so at least one candidate is always
appended.
\end{proof}

\noindent
\textbf{What this does not establish.} The result is a purely
combinatorial fact about prefix sums of a sorted list, true regardless of
how well $\hat p_w$ approximates $p_w$. It says nothing about whether the
\emph{true} top-1 pattern lies in the returned prefix. A statement of
that form would require a finite-sample concentration inequality relating
$\hat p_w$ to $p_w$ --- a multiplicative Chernoff or
Dvoretzky--Kiefer--Wolfowitz-type bound --- yielding something like
$\Pr[\text{true top-1}\notin\text{selected set}]\le g(\dcov,S_w)$. No
such bound is derived here or implemented anywhere in the released code.
$\dcov$ controls coverage of the \emph{sample}, not a calibrated error
rate on the population; the two coincide only in the $S_w\to\infty$
limit. Reading $\dcov$ as an error target, which its name invites, is
therefore incorrect.

\subsection{The carry-aware compatibility check is an exact classifier (Lemma 1)}
\INVAR{}\quad\textit{For overlap width $o>0$, the check accepts iff
$\mathrm{tail}=\mathrm{head}$ or
$(\mathrm{int}(\mathrm{tail},2)-c)\bmod 2^{o}=\mathrm{int}(\mathrm{head},2)$,
where $c$ is the carry bit just past the shared overlap.}

\begin{proof}
The check computes the overlapping bit strings $\mathrm{tail}$ and
$\mathrm{head}$, returns true immediately on exact equality, and
otherwise returns the truth value of exactly the stated modular equation
with $c$ the carry bit just past the shared overlap. There is no other
branch. The ground-truth oracle used by the corrected audit --- ``no
borrow, or exactly one legitimate single-bit borrow'' --- is defined by
the same predicate on the same triple. The two are therefore the same
Boolean function of $(\mathrm{tail},\mathrm{head},c)$ by construction,
not merely in agreement on tested inputs, and the classifier cannot
disagree with the oracle on any input at any overlap width.
\end{proof}

\noindent
Because the equality is algebraic rather than observed, the conclusion
generalises beyond the 168 audited combinations to every
$(\mathrm{tail},\mathrm{head},c)$ triple at any overlap width --- which an
exhaustive audit over a finite input space could not by itself establish.
The label is \INVAR{} rather than \PROVEN{} because the object
constrained is the shipped predicate, not a mathematical abstraction. The
lemma says nothing about whether the \emph{true} bit patterns ever reach
the check; that depends on candidate generation and shot sufficiency.

\subsection{Module C never narrows the beam (Proposition 6)}
\PROVEN{}\quad\textit{Whenever $\beammax\ge P$, we have $P'\ge P$.}

\begin{proof}
$\mathrm{expanded}\ge0$, so $P(1+\mathrm{expanded})\ge P$. Under the
hypothesis $\beammax\ge P$,
$\min(\beammax,P(1+\mathrm{expanded}))\ge\min(\beammax,P)=P$.
\end{proof}

\noindent
The hypothesis holds in every configuration evaluated
($\beammax=4096\gg P\le2$ on the validation grid). Combined with
Proposition 7 below, the consequence is that module C can only
\emph{retain} more candidate paths than baseline at each stitching level,
never fewer. \textbf{What this does not establish:} $\mathrm{expanded}$
is a count of ambiguous windows, not a bound on the number of
globally consistent paths that exist, which depends on the joint
structure of all $W$ windows' candidate sets and the carry check. $P'$ is
therefore a heuristic budget-scaling rule, and nothing proves that the
true path survives pruning at $P'$ on any particular trial.

\subsection{Beam stitching is exact only when unsaturated (Proposition 7)}
\PROVEN{}\quad\textit{If the number of compatible next-paths generated
before truncation never exceeds the beam width at any level, the search
returns the true global optimum under the product-of-weights score.
Conversely, once the beam saturates, discarded paths are permanently
unrecoverable.}

\begin{proof}
If the number of compatible next-paths never exceeds the beam width, the
truncation to the top $P$ is a no-op at every level, so the surviving set
after each level is the full set of compatible paths. By induction over
levels, the final list contains every globally consistent path, sorted by
score, with the true maximum first. Conversely, if the count exceeds the
beam width at some level, the discarded paths are unrecoverable: the loop
only ever extends the surviving set and never re-widens it after
truncation, and there is no backtracking branch.
\end{proof}

\subsection{Module D terminates (Lemma 2)}
\PROVEN{}\quad\textit{The resampling loop for window $w$ halts after at
most $R_w\le\lceil\log_2(\Smax/S_w^{(0)})\rceil$ rounds.}

\begin{proof}
Each round adds $\min(S_w,\Smax-S_w)$ shots, where $S_w$ is the current
total. While $S_w\le\Smax/2$ the added amount is $S_w$ itself, so $S_w$
at least doubles, giving at most $\lceil\log_2(\Smax/S_w^{(0)})\rceil$
such rounds before $S_w>\Smax/2$. Thereafter the added amount is
$\Smax-S_w<S_w$, and $S_w$ is strictly increasing and bounded above by
$\Smax$, so the loop reaches $S_w=\Smax$ --- at which point the increment
is non-positive and the loop exits --- within one further round.
\end{proof}

\noindent
This is a genuine termination guarantee rather than a mere engineering
cap, and it matches the configured defaults exactly: at $S_w^{(0)}=4$ and
$\Smax=64$, $\lceil\log_2(64/4)\rceil=4$ rounds, which is the value the
configuration file's own comment records.

\subsection{Equivalence to the published baseline (Invariant 1)}
\INVAR{}\quad\textit{With all three modules disabled, the adaptive entry
point executes a call sequence structurally identical to the published
pipeline.}

\begin{proof}[Justification (code inspection, not a mathematical proof)]
By exhaustive inspection of the adaptive entry point. Each of the three
module flags guards a complete \texttt{if}/\texttt{else}, and every
\texttt{else} branch reproduces the pre-adaptive computation exactly:
the same function calls with the same arguments, sourced from the fixed
\texttt{candidate\_count} and \texttt{max\_paths} configuration fields
rather than from any adaptive quantity. With all three flags false, the
shot-stopping call is skipped entirely by its guard, so counts are
sampled once and unmodified. No adaptive code path is reachable.
\end{proof}

\noindent
The argument is complete for the implementation as written and carries no
force for a re-implementation. Two independent execution-level checks
accompany it, both reported in the main article: the Baseline arm's
success rate reproduces the original campaign's figure across three
measurements (52.1\%, 52.1\%, 51.1\%), where the original campaign drove
a \emph{different} entry point, making the agreement cross-implementation
evidence; and an instrumented re-run of the canonical grid reproduced the
recorded per-trial outcomes of all five arms on all 144 trials
bit-for-bit.

\subsection{Sufficient decomposition of end-to-end correctness}
\PROVEN{}\quad\textit{If (i) every window's top-ranked retained candidate
equals the true window bit pattern, (ii) the carry check accepts exactly
the carry-consistent pairs among those patterns, (iii) beam pruning never
discards the resulting fully-true path, and (iv) continued-fraction
expansion recovers the true order, then the pipeline returns the correct
factorisation.}

\begin{proof}
Under (i)--(ii) the true patterns are pairwise compatible across every
overlap, so a fully-true path exists at every stitching level. Under
(iii) it survives pruning. Since the true pattern is top-ranked in every
window by (i), that path attains the maximal product score and is ranked
first. Under (iv), order recovery on it succeeds, so the outer loop
returns on its first iteration.
\end{proof}

\noindent
This decomposition is a logical convenience for organising the results
above, \textbf{not} a claim that (i)--(iv) are \emph{necessary}. A correct
answer can also arise from a non-top-ranked path, from a legitimate
carry-borrow substituting for exact equality, or from a later stitched
candidate succeeding where an earlier one failed. The sufficiency framing
is retained because it is what the component-level results can actually
establish.

\subsection{Monotone non-restrictiveness, and why it does not imply monotone outcomes}
\PROVEN{}\quad Under the per-module hypotheses above, each of B, C and D
can only \emph{relax} a resource constraint relative to baseline --- more
candidates, a wider beam, more shots --- and never tighten one. No module
removes information or candidates that baseline would have kept.

It is tempting to read this as implying that no arm can underperform
baseline on any trial. That inference does not hold, and the data refute
it: on the canonical grid module B alone turned two baseline successes
into failures and module D six, and on $(N{=}15,a{=}2,r{=}4)$ module D
broke exactly as many trials as it fixed. The reason is that relaxing a
budget changes which candidates are retained and how the resulting paths
rank, and continued-fraction order recovery is not monotone in the
candidate set: a different phase can be tried first and fail where the
baseline's first choice would have succeeded. \textbf{The monotonicity is
a statement about resources, not about outcomes}, and the aggregate gains
reported in the main article are empirical results rather than corollaries
of this structure.

% =====================================================================

% =====================================================================
\section{Calibration: bin-level data and diagnostics}
\label{si:calib}

Table~\ref{tab:calib} is the numerical form of the calibration figure in
the main article, over 2{,}687 per-window observations from 579
recaptured trials. Bins below $0.4$ are empty by construction:
$\Cw\ge0.5$ whenever $n_1\ge n_2$, and $b^{(1)}_w$ is the top-ranked
pattern by definition, so values below $0.5$ arise only from ties
resolved by ordering.

\begin{table}[h]
\centering\small
\caption{Reliability of $\Cw$ by predicted-confidence bin, from
\texttt{phase6\_calibration.csv}. Raw Brier score $0.2239$ against
$0.2500$ for a constant always-$0.5$ forecast; isotonic recalibration
reduces this to $0.1988$ in-sample.}
\label{tab:calib}
\begin{tabular}{@{}lccc@{}}
\toprule
Predicted $\Cw$ bin & $n$ windows & Mean predicted & Observed correct \\
\midrule
$(0.4, 0.5]$ & 340 & 0.500 & 0.512 \\
$(0.5, 0.6]$ & 216 & 0.518 & 0.606 \\
$(0.6, 0.7]$ & 227 & 0.657 & 0.537 \\
$(0.7, 0.8]$ & 743 & 0.749 & 0.693 \\
$(0.8, 0.9]$ & 280 & 0.859 & 0.639 \\
$(0.9, 1.0]$ & 881 & 0.964 & 0.792 \\
\midrule
All          & 2{,}687 & --- & --- \\
\bottomrule
\end{tabular}
\end{table}

\paragraph{The overconfidence is not uniform across the range.} The
lowest bin is well calibrated ($0.500$ predicted against $0.512$
observed), and the $(0.5,0.6]$ bin is if anything under-confident. The
gap opens in the upper half and is widest in $(0.8,0.9]$ --- $0.859$
predicted against $0.639$ observed --- and in the heavily populated top
bin, $0.964$ predicted against $0.792$ observed. Since module D's
stopping threshold at $\varepsilon=0.05$ sits inside that top bin, the
practical consequence is the one stated in the main article: windows that
halt resampling are wrong about one time in five, not one in twenty.

\paragraph{Ground truth.} The reference is not a bit-slice of a single
``true phase''. The standard Shor $|1\rangle$ initial state is an equal
superposition over all $r$ eigenphases $s/r$, so no single reference bit
pattern per window exists. The well-defined ground truth is instead
whether the observed top-1 pattern's \emph{true} (noiseless,
infinite-shot) marginal probability really does exceed the observed
top-2 pattern's --- that is, whether $\Cw$'s claim is correct for that
window. The true marginal distribution is computed once per
$(a,N,\text{window spec})$ by exact statevector simulation, independent
of shots and seed, and cached.

\paragraph{Mechanism.} The miscalibration is the predicted consequence of
assumption (A2) being false. $b^{(1)}_w$ and $b^{(2)}_w$'s counts are two
negatively correlated coordinates of one $K_w$-way multinomial, not two
independent Binomials; ignoring that correlation pushes $\Cw$ toward more
extreme values than the joint comparison would give. The defect is
therefore structural rather than a tuning failure, and it is monotone
rather than erratic --- which is precisely why the raw statistic remains
serviceable as a ranking signal.

\paragraph{Recalibration, and why it is not shipped.} Isotonic regression
\cite{zadrozny2002} maps the raw score to a Brier of $0.1988$, bringing
the top-bin value to $0.82$ against the same $0.792$ observed. This is an
\emph{in-sample} fit on the same 2{,}687 observations, with no held-out
split and no cross-validation, so it bounds what recalibration could
achieve rather than estimating what it would achieve out of sample. It is
deliberately not wired into module D's stopping rule, for two reasons.
First, doing so would make the shipped rule depend on a fit whose
generalisation to other instances, orders and window geometries has not
been tested. Second, it would break the correspondence between the
shipped configuration and the reported ablation, which thresholds on raw
$\Cw$ throughout. Substituting the recalibrated value is a reasonable
follow-up, but it would require re-running the ablation to attribute any
resulting change correctly.

% =====================================================================

% =====================================================================
\section{Artifact map and reproducibility detail}
\label{si:artifacts}

\subsection{Script-to-result-to-section mapping}

Every numeric claim in the main article and in this document is a
mechanical summary of one of the files in Table~\ref{tab:artifacts}. All
files are in the released repository under
\texttt{Data/failure\_study/}, and all generating scripts under
\texttt{Experiments/}. The table lists only the files this article
consumes; the companion study~\cite{companion} lists the files supporting
the end-to-end evaluation, from the same release.

\begin{table}[h]
\centering\small
\caption{Artifact map for the method article. ``Used in'' names the
section of the main article (M) or of this document (S) that consumes the
file.}
\label{tab:artifacts}
\scriptsize
\begin{tabular}{@{}>{\raggedright\arraybackslash}p{0.31\linewidth}>{\raggedright\arraybackslash}p{0.30\linewidth}>{\raggedright\arraybackslash}p{0.31\linewidth}@{}}
\toprule
Result file & Produced by & Used in \\
\midrule
\texttt{raw\_\allowbreak{}windows/*.json} (579 files) & \texttt{recapture\_\allowbreak{}window\_\allowbreak{}counts.py} & M: method validation \\
\texttt{paperA\_\allowbreak{}window\_\allowbreak{}validation.csv} & \texttt{paperA\_\allowbreak{}window\_\allowbreak{}validation.py} & M: Table~5, Fig.~2 \\
\texttt{phase6\_\allowbreak{}calibration.csv} & \texttt{phase6\_\allowbreak{}calibration.py} & M: Fig.~3; S: Table~\ref{tab:calib} \\
\midrule
\texttt{plots/paper\_\allowbreak{}figA\_\allowbreak{}window\_\allowbreak{}retention.png} & \texttt{paperA\_\allowbreak{}window\_\allowbreak{}validation.py} (rendering only) & M: Fig.~2 \\
\texttt{plots/paper\_\allowbreak{}fig5\_\allowbreak{}calibration.png} & \texttt{paper\_\allowbreak{}figures.py} (rendering only; runs no trial) & M: Fig.~3 \\
\bottomrule
\end{tabular}
\end{table}

\subsection{Determinism and the analysis pipeline}

\paragraph{Trials are deterministic given their seed.} Each trial's seed
is derived from its grid coordinates and seeds the simulator's sampling,
so re-running any cell reproduces its counts exactly. It was verified
rather than assumed: an instrumented re-run reproduced all 144 canonical
trials, in all five arms, bit-for-bit~\cite{companion}.

\paragraph{The window-level analysis samples nothing.} The retention
script replays the frozen per-window counts in
\texttt{raw\_windows/} through the shipped candidate-generation
functions and scores the retained sets against statevector marginals
computed once per $(a,N,\text{window specification})$. It executes no
circuit on the sampler and recomputes no recorded outcome, and it asserts
the calibration corpus size and raw Brier score before writing any
output.

\paragraph{The reference implementation carries its own checks.} Every
executable module runs a self-test without arguments; the adaptive
module's asserts Invariant~1 directly, by comparing the disabled-module
path against the published fixed-budget reconstruction on identical
counts.

\bibliographystyle{elsarticle-num}
\bibliography{refs_A}

\end{document}
